\par\clearpage\NoSourceHead\SourceLeaf{296}
\begin{center}\textbf{Sachregister.}\end{center}
\begingroup
\newcommand{\SI}[1]{\par\noindent\hangindent=1.3em\hangafter=1 #1\par}
\newcommand{\SG}{\par\vspace{.7\baselineskip}}
\setlength{\tabcolsep}{4pt}
\noindent\begin{tabular}{@{}p{.48\linewidth}|p{.48\linewidth}@{}}
\begin{minipage}[t]{\linewidth}\fontsize{8.2}{10.5}\selectfont\raggedright
\SI{Abbildung zweier Flächen aufeinander 172ff.}
\SI{Absolute Eigenschaften einer Fläche 175}
\SI{Absolute Geometrie 178}
\SI{Absolute Invariante der Modulgruppe 101}
\SI{Abstandsbegriff 50}
\SI{Additionstheorem 32, 73ff., 168}
\SI{Ähnliche und verkehrte Lage 174, 179ff.}
\SI{Al-Funktionen, \textsc{Weierstrasssche}, im lemniskatischen Falle 35}
\SI{Algebraische Theoreme 189ff.}
\SI{Algebraische Theorie des agM. 106}
\SI{Analytische Fortsetzung 153, 154, 186}
\SI{Arithmetisch-geometrisches Mittel 11ff., 18, 24, 27, 29, 39, 40, 41ff., 62ff., 70ff., 85ff., 89ff., 94, 106, 112, 129, 140, 145, 150, 152, 166ff., 184, 198}
\SI{— Zusammenhang mit dem Ellipsenquadranten 65, 66, 166ff.}
\SI{— — mit dem vollständigen Integral erster Gattung 69, 89, 94, 169}
\SI{Asymptotisches Verhalten 24, 26, 50, 58, 60, 67, 115ff.}
\SI{Attraktion der Sphäroide 154, 155}
\SI{Automorphe Funktionen 198}
\SI{Automorphismus der Modulfunktion 87}
\SG
\SI{Bahnbestimmung kleiner Planeten 118ff.}
\SI{Benachbarte quadratische Formen 99}
\SI{\textsc{Besselsche} Funktionen 26, 50}
\SI{Bilinearer Algorithmus 130ff.}
\SI{Biquadratische Reste 52, 124, 188}
\SI{Briefwechsel, gedruckte 4}
\SG
\SI{\textsc{Cauchy-Riemannsche} partielle Differentialgleichungen 172}
\SG
\SI{Darstellung einer Zahl durch vier Quadrate 83, 182}
\SI{\textsc{Descartessche} Zeichenregel siehe \textsc{Harriotscher} Satz.}
\SI{Differentialgleichung der elliptischen Funktionen 73, 80ff., 97}
\end{minipage}&
\begin{minipage}[t]{\linewidth}\fontsize{8.2}{10.5}\selectfont\raggedright
\SI{Differentialrelationen für Thetafunktionen 83}
\SI{Dimensionen des Raumes 179ff.}
\SI{Dioptrische Untersuchungen 199}
\SI{\textsc{Dirichletsches} Prinzip 194}
\SI{Diskriminante (Determinante) 158}
\SI{Doppelte Periodizität}
\SI{— im lemniskatischen Falle 33, 51}
\SI{— allgemein 74}
\SI{Drehungen einer Kugel 105, 198}
\SG
\SI{Einfachstes agM. 104}
\SI{\textsc{Eisensteinsche} Reihen 35}
\SI{Elementare Mittel 44}
\SI{Ellipsenquadrant 64, 65, 66, 166ff.}
\SI{Elliptische Funktionen 28, 40, 41, 52, 71, 72ff., 85, 127, 129, 132, 145, 183, 198}
\SI{— Funktion Sl 75ff., 85, 97, 108, 110, 186, 198}
\SI{— — \textsc{Jacobische} 85, 186, 193, 197}
\SI{Elliptische Integrale 27ff.\hfill [186}
\SI{— Integral erster Gattung 70ff., 108, 110, 132, 152,}
\SI{— — vollständiges 69ff., 138, 152, 169}
\SI{— — zweiter Gattung siehe Rektifikation der Ellipse und Ellipsenquadrant}
\SI{Empirischer Charakter des Raumes 177}
\SI{Entwicklung nach trigonometrischen Funktionen}
\SI{— für sin lemn 38}
\SI{— für Zähler und Nenner von sin lemn 38, 39}
\SI{— diskontinuierlicher Funktionen 199}
\SI{— periodischer Funktionen 202ff.}
\SI{Erdmagnetismus 188}
\SI{\textsc{Eulersche} Integrale 140, 146, 147, 153}
\SI{Exhaustionsmethode 48}
\SG
\SI{Fakultäten, \textsc{Krampsche}, 135}
\SI{Frequenz der Primzahlen 135}
\SI{Funktionalgleichung für \(P(n,x)\) 115}
\SI{— für \(\Pi(z)\) 146}
\SI{Funktionentheorie, allgemeine 113, 148, 151ff., 186, 205}
\end{minipage}
\end{tabular}
\HistoricalFoot{}{27\textsuperscript{*}}
\clearpage\SourceLeaf{297}\ArithmeticHead{212}{L}{L. SCHLESINGER, ÜBER GAUSS' FUNKTIONENTHEORETISCHE ARBEITEN.}
\noindent\begin{tabular}{@{}p{.48\linewidth}|p{.48\linewidth}@{}}
\begin{minipage}[t]{\linewidth}\fontsize{8.2}{10.5}\selectfont\raggedright
\SI{\textsc{Gausssche} Differentialgleichung 90, 91, 112, 136, 138, 140, 148, 149, 151}
\SI{\textsc{Gausssche} Identität 159}
\SI{\textsc{Gausssche} Reihe \(F(\alpha,\beta,\gamma,x)\) 27, 28, 64, 90, 91, 112, 124, 133, 135, 136ff., 145, 184, 198, 199}
\SI{\textsc{Gaussscher} Satz der Potentialtheorie 163}
\SI{\textsc{Gausssche} Summen 122}
\SI{\textsc{Gausssche} Transformation 78, 108ff., 131, 169}
\SI{Geodäsie 170, 199}
\SI{Geometrische Funktionentheorie 111, 191, 206}
\SI{\textsc{Greensche} Funktion 193}
\SI{Grenzwert (Limes, absolute Grenze) 49, 50, 55}
\SI{— für das agM. 89}
\SI{Grundlagen der Analysis 44ff.}
\SI{Grundlagen der Geometrie 3, 47, 63, 209ff.}
\SI{Gruppe, ganzzahliger linearer Transformationen (Modulgruppe) 95}
\SG
\SI{Harmonisch-arithmetisches Mittel 24, 44}
\SI{Harmonisch-geometrisches Mittel 24, 87}
\SI{\textsc{Harriotscher} Satz 187}
\SI{Hauptfeld 105}
\SI{Hauptkongruenzgruppe 2. Stufe 98, 103}
\SI{— — ihr Fundamentalbereich 103, 104, 105, 111, 113}
\SI{— 4. Stufe 99}
\SI{— — ihr Fundamentalbereich 104}
\SI{\textsc{Hermitesches} Transformationsprinzip 77ff.}
\SI{Hypergeometrische Reihen 58}
\SG
\SI{Imaginär siehe Komplex}
\SI{Inexplikable Funktionen \(\Pi\) und \(\Psi\) 115, 139, 143, 146ff.}
\SI{Integrallogarithmus 134, 135, 152}
\SI{Integraltheoreme der Potentialtheorie 155, 188, 209}
\SI{Integration im komplexen Gebiete 116, 152, 153, 162, 165, 174, 203ff.}
\SI{Interpolation 124ff.}
\SI{Isometrische Abbildung (Abwicklung) 175}
\SG
\SI{Kettenbruch 23, 24, 114, 138}
\SI{Kettenbruchentwicklung für den Quotienten \textsc{Gaussscher} Reihen 137, 140, 146, 156}
\SI{Kettenbruchentwicklung für die Substitutionen der Modulgruppe 101, 102}
\SI{— der Hauptkongruenzgruppe 2. Stufe 98}
\SI{Klassenanzahl quadratischer Formen 116ff.}
\SI{Klassenbegriff für lineare Differentialgleichungen 146}
\SI{Kleine Planeten 117ff., 125}
\SI{Koeffizienten von \textsc{Laplace} 27, 138, 144}
\end{minipage}&
\begin{minipage}[t]{\linewidth}\fontsize{8.2}{10.5}\selectfont\raggedright
\SI{Komplexe Veränderliche 34, 51, 52ff., 71, 147, 171}
\SI{— Zahlen oder Grössen 51, 57, 60, 147, 191, 203}
\SI{— — ihre Metaphysik 53, 176}
\SI{— — ihre geometrische Darstellung 53ff., 60}
\SI{Konforme Abbildung 170ff., 191, 205}
\SI{— — der Ellipse 192ff.}
\SI{Konvergenz und Divergenz von Reihen 50, 58, 140ff.,}
\SI{— der Thetareihen 79ff., 86\hfill [202ff.}
\SI{— von \(F(\alpha,\beta,\gamma,x)\) 144}
\SI{— von \(F(\alpha,\beta,\gamma,1)\) 143, 146}
\SI{Konvergenzuntersuchung}
\SI{— für lemniskatische Funktionen 35}
\SI{— für das \(\Pi(z)\) definierende unendliche Produkt 146}
\SI{— für die \(\Pi\), \(\Pi\Psi\) und \(\Psi\) darstellenden Integrale 147}
\SI{Kreisteilung 20ff., 51, 52, 127}
\SI{Kugelfunktionen 156}
\SG
\SI{\textsc{Lagrangescher} Umkehrungssatz 25, 26, 66, 199}
\SI{\textsc{Landensche} Transformation 12, 29, 78, 108ff., 130, 168}
\SI{\textsc{Laplacesche} partielle Differentialgleichung 156, 186}
\SI{\textsc{Leist}notizen, reguläre, irreguläre 7}
\SI{Lemniskatische Funktionen 28, 30ff., 51, 73ff., 123, 124, 129, 168, 198}
\SI{— — sin lemn, cos lemn 33, 38, 39, 124}
\SI{Lemniskatische Integrale 29, 52, 72, 73, 127, 168}
\SI{— vollständige 31ff.}
\SI{Lemniskatische Periode 33ff., 43, 51, 66, 146, 185}
\SI{Lineare Differentialgleichung, allgemeine Theorie 151}
\SI{— —, zweiter Ordnung 112, 143, 198\newline\hspace*{1.3em}siehe auch \textsc{Gausssche} Differentialgleichung}
\SI{Lineare Punktmengen, beliebige 49}
\SI{Logarithmentafeln}
\SI{—, \textsc{Gausssche} 155, 201}
\SI{—, \textsc{Schulzesche} 11, 12, 114}
\SI{—, \textsc{Wolframsche} 114}
\SI{Logarithmus 149, 152, 164}
\SG
\SI{Mathematik, niedere und höhere 48}
\SI{Mechanische Quadratur 156}
\SI{Mehrdeutigkeit des agM. 86, 93ff., 99, 100}
\SI{Metaphysik der Mathematik 45ff.}
\SI{— — Raumlehre 176, 181}
\SI{— — negativen und imaginären Grössen 176}
\SI{Methode der kleinsten Quadrate 12}
\SI{— —, Vorlesung 206, 207ff.}
\SI{Minimalkurven einer Fläche 173}
\SI{Mittelpunktsgleichung 50, 115ff., 125, 141, 200, 202}
\SI{Mittelwertsatz der Potentialtheorie 163}
\SI{Mittlere Werte der Zahlentheorie 11, 24, 26}
\end{minipage}
\end{tabular}
\clearpage\SourceLeaf{298}\ArithmeticHead{213}{R}{SACHREGISTER.}
\noindent\begin{tabular}{@{}p{.48\linewidth}|p{.48\linewidth}@{}}
\begin{minipage}[t]{\linewidth}\fontsize{8.2}{10.5}\selectfont\raggedright
\SI{Modulfunktion 87ff., 93ff., 98, 99, 111, 113, 124, 145, 184}
\SI{—, Aufzeichnungen zur Theorie der, 93, 94}
\SI{Modulgruppe 101, 198}
\SI{— ihr Fundamentalbereich 105, 198}
\SI{Monogene Funktion des Ortes auf einer Fläche 174ff.}
\SI{Multiplikationsformeln elliptischer Funktionen 197}
\SI{— im lemniskatischen Falle 32, 33, 52}
\SG
\SI{Nachlassteile, Liste der benutzten, 8}
\SI{Natürliche Grenze 106}
\SI{Nichteuklidische Geometrie 175ff.}
\SI{Null- und Unendlichkeitsstellen der lemniskatischen Funktionen 35}
\SI{Numerisch-exponentieller Kalkul 114}
\SG
\SI{Parallelenaxiom 23, 24, 47}
\SI{Pentagramma mirificum 195ff.}
\SI{Potential}
\SI{—, ebenes (logarithmisches) 156, 186, 191, 192, 200,}
\SI{—, räumliches 156, 187, 193\hfill [201, 209}
\SI{Potentialtheorie 156, 205}
\SI{Potenzreihen 15, 50, 58, 66}
\SI{—, deren Exponenten eine Progression 2. Ordnung bilden, siehe Summatorische Funktionen}
\SI{Potenzreste, höhere 123}
\SI{Produktdarstellung der uniformisierenden Variablen}
\SI{Punktgitter, quadratisches 34, 117\hfill [des agM. 107}
\SG
\SI{Quadratische Formen, binäre 88, 95, 98ff., 102, 116}
\SI{Quellen, Verzeichnis der benutzten 4}
\SG
\SI{Randwertproblem 193, 194}
\SI{Realitätsverhältnisse 70ff.}
\SI{Rechts und links 55, 171ff.}
\SI{Reduktionstheorie der quadrat. Formen 98ff., 198}
\SI{Reelle Zahl oder Grösse 48}
\SI{Reihen (Folgen) von reellen Grössen 49, 140ff.}
\SI{Rektifikation der Ellipse und Hyperbel 27, 29, 64, 70, 91, 108, 110, 126, 127, 169}
\SI{— — Lemniskate 28, 32, 52, 70}
\SI{Relationen zwischen verwandten Funktionen 64, 137, 139, 146, 148, 149}
\SI{Reziprozitätsgesetz der quadrat. Reste 21, 22, 24}
\SI{Rigor antiquus 48, 59, 141, 206}
\SG
\SI{Schliessungsproblem 195, 196, 197}
\SI{Schöner Lehrsatz 107, 140}
\SI{Schranke (Grenze), obere, untere 49}
\SI{Semikonvergente Reihen 146, 147}
\end{minipage}&
\begin{minipage}[t]{\linewidth}\fontsize{8.2}{10.5}\selectfont\raggedright
\SI{Spiegelung an Kreisen 105}
\SI{Summatorische Funktionen 14ff., 18, 23, 29, 65, 69, 70, 76ff., 93, 94ff., 182}
\SI{— im lemniskatischen Falle 39}
\SI{Summe divergenter Reihen 50}
\SI{Symmetrische Funktionen 158}
\SI{Symmetrische Körper 179ff.}
\SI{Störungsfunktion 26ff., 89, 90}
\SG
\SI{Tagebuch 5, 20, 118, 154}
\SI{Teilungsgleichung}
\SI{— für die lemniskatischen Funktionen 33, 52}
\SI{— — Fünf- und Siebenteilung 36, 37}
\SI{— für die Perioden 84, 129}
\SI{Theorema elegantissimum 107}
\SI{Thetafunktionen}
\SI{— allgemein 75ff., 95ff., 128, 182, 198}
\SI{— Identität zwischen Summen- und Produktdarstellung 16, 82, 123, 129, 182}
\SI{— im lemniskatischen Falle 38, 168}
\SI{— Produktdarstellung 82}
\SI{— zweiter Ordnung 77ff., 81ff.\newline\hspace*{1.3em}siehe auch Summatorische Funktionen}
\SI{Topologie schlichter ebener Bereiche 203ff.}
\SI{Topologische Abbildung 171}
\SI{Transformation der elliptischen Funktionen 186}
\SI{— lineare der Thetafunktionen 96, 128}
\SI{— — der summatorischen Funktionen 96, 101, 122}
\SI{— zweiter Ordnung der Thetafunktionen 77ff., 83, 185}
\SI{Transzendente Theorie des agM. 106}
\SI{Trinomische Gleichungen 200, 201ff.}
\SG
\SI{Umkehrungsfunktion von Integralen 30, 74}
\SI{— des allgem. elliptischen Integrals erster Gattung 74, 85, 107, 110}
\SI{— des Integrals \(\displaystyle\int\frac{dx}{\sqrt{1\pm x^3}}\) 36, 73}
\SI{— des lemniskat. Integrals 32}
\SI{Umlauf (reditus) 151, 205}
\SI{Unbestimmtheitsgrenzen (lim sup, lim inf) 49}
\SI{Uniformisierende Variable für das agM. 68, 69}
\SI{— rechnerische Methoden zu ihrer Bestimmung 92,}
\SI{Uniformisierung\hfill [107}
\SI{— des agM. erste, unzureichende 42}
\SI{— — durch die summatorischen Funktionen 95}
\SI{— des Ellipsenquadranten 65, 66}
\SG
\SI{Valenz 101}
\SI{Vektoren in der Ebene 54}
\SI{Vierter Einwand gegen \textsc{D’Alembert} 58ff., 162, 194}
\end{minipage}
\end{tabular}
\clearpage\SourceLeaf{299}\ArithmeticHead{214}{L}{L. SCHLESINGER, ÜBER GAUSS' FUNKTIONENTHEORETISCHE ARBEITEN.}
\noindent\begin{tabular}{@{}p{.48\linewidth}|p{.48\linewidth}@{}}
\begin{minipage}[t]{\linewidth}\fontsize{8.2}{10.5}\selectfont\raggedright
\SI{»Wahre« Geometrie 180ff.}
\SI{Wahrscheinlichkeitsrechnung 114}
\SI{Widerspruchsfreiheit der nichteuklidischen (absoluten) Geometrie 176, 180}
\SI{Witwenkasse 199, 202}
\SI{Wurzelexistenz, Beweis der 24, 58, 59}
\SI{— erster (Inauguraldissertation) 45, 53, 57ff., 111, 165, 190, 191, 200}
\end{minipage}&
\begin{minipage}[t]{\linewidth}\fontsize{8.2}{10.5}\selectfont\raggedright
\SI{Wurzelexistenz, Beweis, zweiter 59, 154, 157ff.}
\SI{— dritter 59, 157, 161ff.}
\SI{— vierter 60, 200ff.}
\SI{Wurzelzeichen, Lösung von algebraischen Gleichungen durch, 59}
\SG
\SI{Zusammengehörige Werte des agM. 80, 96, 97}
\SI{Zusammenhang ebener Bereiche 174}
\end{minipage}
\end{tabular}
\par\bigskip\begin{center}\rule{.19\linewidth}{.3pt}\end{center}
\endgroup
\ArithmeticHead{214}{L}{L. SCHLESINGER, ÜBER GAUSS' FUNKTIONENTHEORETISCHE ARBEITEN.}
